\documentclass[10pt,a4paper]{article}
\usepackage[utf8]{inputenc}
\usepackage[T1]{fontenc}
\usepackage{amsmath,amssymb}
\usepackage{graphicx}
\usepackage{booktabs}
\usepackage{hyperref}
\usepackage{xcolor}
\usepackage[margin=2.5cm]{geometry}

\title{Chim\`ere: A Self-Improving MoE Inference System\\for Consumer GPU Hardware}

\author{
  K\'evin Remondi\`ere\\
  Independent Researcher, Oloron-Sainte-Marie, France\\
  \texttt{kevin.remondiere@gmail.com}\\
  ORCID: \href{https://orcid.org/0009-0008-2443-7166}{0009-0008-2443-7166}\\
  \url{https://github.com/AIdevsmartdata/chimere}
}

\date{March 2026}

\begin{document}
\maketitle

\begin{abstract}
We present Chim\`ere, an integrated system for running a 35B-parameter Mixture-of-Experts language model (Qwen3.5-35B-A3B) on a single RTX~5060~Ti (16\,GB VRAM) at 93~tokens/second, with continuous self-improvement through a nightly quality loop. The system combines five subsystems: (1)~a Rust-native inference runtime with custom CUDA~sm\_120 kernels (56K~lines), (2)~multi-tier Engram memory with domain-specific N-gram tables, (3)~an intent-classifying orchestrator with entropy-adaptive compute routing, (4)~a SOTA web search pipeline with 8-stage retrieval, and (5)~a quality-gated training data pipeline feeding nightly LoRA fine-tuning. On a 10-question benchmark spanning physiotherapy, code generation, mathematics, and tool calling, the integrated system achieves 100\% pass rate (vs.\ 50\% before integration), with Engram v2 ($\alpha{=}0.1$, response-only) providing 88\% accuracy (vs.\ 85\% without, 77\% with naive Engram~v1). We release the complete system as open source across three repositories totaling 121K~lines of code.
\end{abstract}

\section{Introduction}

The gap between frontier API models (GPT-4o, Claude~Opus) and local deployment remains significant. Consumer GPUs with 16\,GB VRAM can run quantized 7--13B models comfortably, but 35B+ models require aggressive quantization that degrades quality. The standard approach --- quantize and serve --- ignores the system-level optimizations possible when the model, orchestration, retrieval, and training loop are co-designed.

Chim\`ere is a \emph{system}, not a model. It combines:
\begin{itemize}
  \item A \textbf{Rust inference runtime} (chim\`ere-server) with custom CUDA kernels optimized for NVIDIA Blackwell (sm\_120), achieving 93~tok/s on Qwen3.5-35B-A3B at IQ3\_S quantization (3.56~BPW, 14.71~GB).
  \item An \textbf{orchestrator} (ODO) that classifies intent, enriches context with RAG and web search, and routes requests to appropriate compute profiles.
  \item A \textbf{self-improvement loop} that logs quality scores from production traffic, generates training pairs, and runs nightly LoRA fine-tuning and Engram updates.
\end{itemize}

The key insight is that these components are synergistic: the orchestrator's quality scores improve the training data, the training improves the model, and the improved model generates better quality scores. The system improves while idle.

\section{System Architecture}

\begin{verbatim}
User request --> ODO (port 8084)
                  |-- Intent classifier (regex->filetype->LLM)
                  |-- Enricher (web search, RAG, tools)
                  |-- Entropy router (fast/quality/ultra)
                  |-- Engram v2 (alpha=0.1, response-only)
                  |-- DVTS tree search (ThinkPRM scoring)
                --> chimere-server (port 8081)
                  |-- Rust FFI to ik_llama (IQ3_S GEMV)
                  |-- Multi-tier Engram lookup
                  |-- Logprobs (top-5 log-softmax)
                --> ThinkPRM (port 8085, CPU)
                --> quality_scores.jsonl
                --> nightly: Engram WRITE + LoRA + DSPy
\end{verbatim}

\section{Inference Runtime}

\subsection{Hardware-Optimized Serving}

Chim\`ere-server is a 56{,}453-line Rust application with 2{,}588 lines of custom CUDA kernels targeting NVIDIA Blackwell (sm\_120). It provides an OpenAI-compatible API and integrates with \texttt{ik\_llama.cpp} via FFI for quantized matrix operations.

\begin{table}[h]
\centering
\caption{Inference performance on RTX~5060~Ti 16\,GB.}
\begin{tabular}{lcc}
\toprule
\textbf{Configuration} & \textbf{Gen (tok/s)} & \textbf{VRAM} \\
\midrule
IQ3\_S custom-mix, ik\_llama sm120, ncmoe=2, 10K ctx & 98.6 & 15.2\,GB \\
IQ3\_S custom-mix, ik\_llama sm120, ncmoe=0, 64K ctx & 91.0 & 14.5\,GB \\
IQ3\_S custom-mix, stock llama.cpp, ncmoe=4, 128K ctx & 72.1 & 15.2\,GB \\
IQ3\_S Unsloth UD, stock llama.cpp, 3$\times$64K ctx & 69.0 & 15.8\,GB \\
Q4\_K\_M, ik\_llama sm120, ncmoe=14, 64K ctx & 62.3 & 15.1\,GB \\
Q5\_K\_XL, ik\_llama sm120, ncmoe=18, 64K ctx & 49.9 & 15.0\,GB \\
\bottomrule
\end{tabular}
\label{tab:perf}
\end{table}

The ik\_llama sm\_120 build provides +23\% generation throughput over stock llama.cpp, measured consistently across IQ3\_S ($91 \to 69$~tok/s), Q4\_K\_M ($62 \to 52$), and dense 9B ($44 \to 41$).

\subsection{RAMP Quantization}

We developed RAMP (RL-guided Adaptive Mixed-Precision), a data-free quantization pipeline that assigns per-tensor precision based on NSDS sensitivity analysis and evolutionary search. RAMP~v2 produces a 15.2\,GB GGUF (3.78~BPW) with IQ3\_S base and Q8\_0/Q6\_K/Q5\_K overrides for critical paths (SSM gates, attention QKV, shared expert). The RAMP pipeline and model are released separately.\footnote{\url{https://github.com/AIdevsmartdata/ramp-quant}}

\section{Multi-Tier Engram Memory}

Inspired by the DeepSeek Engram paper~(arXiv:2601.07372), we implement a multi-tier N-gram memory system:

\begin{table}[h]
\centering
\caption{Engram memory tiers.}
\begin{tabular}{llll}
\toprule
\textbf{Tier} & \textbf{Mechanism} & \textbf{Latency} & \textbf{Purpose} \\
\midrule
0 & Cuckoo filter & $<$10\,ns & Existence check (skip 97\% lookups) \\
1 & FNV-1a hash tables & O(1) & Core N-gram lookup \\
2 & FAISS semantic & $\sim$5\,ms & Few-shot example retrieval \\
\bottomrule
\end{tabular}
\end{table}

Domain tables: code (155\,KB), cyber (172\,KB), general (679\,KB). Tables are built by tokenizing domain documents, extracting N-grams, and storing in a binary hash table format read by both the Rust runtime and Python tools.

\subsection{Engram Ablation}

We measured Engram impact on our 10-question benchmark:

\begin{table}[h]
\centering
\caption{Engram ablation results.}
\begin{tabular}{lcc}
\toprule
\textbf{Configuration} & \textbf{Accuracy} & \textbf{Notes} \\
\midrule
Engram v1 ($\alpha{=}0.35$, think+response) & 77\% & Biases thinking phase \\
Engram OFF ($\alpha{=}0$) & 85\% & Baseline \\
\textbf{Engram v2} ($\alpha{=}0.1$, response-only) & \textbf{88\%} & Production \\
\bottomrule
\end{tabular}
\end{table}

The key insight: applying Engram bias during the thinking phase \emph{degrades} performance because it constrains the model's reasoning with domain patterns. Restricting bias to the response phase and lowering $\alpha$ from 0.35 to 0.1 yields a +11~pp improvement over v1 and +3~pp over no Engram.

\section{Intent Classification and Compute Routing}

ODO (1{,}525~lines Python) classifies incoming requests using a 3-strategy cascade:

\begin{enumerate}
  \item \textbf{Keyword/regex matching} ($<$1\,ms, handles $\sim$99\% of requests): pattern-based route assignment to code, kine, cyber, research, or default.
  \item \textbf{File-type detection}: identifies images, code files, data files for specialized routing.
  \item \textbf{LLM fallback with GBNF constraint} ($\sim$50\,ms, $<$1\% of requests): the model classifies the request using a grammar-constrained output.
\end{enumerate}

The entropy router measures request complexity (query length, classifier confidence, historical quality for similar queries) and assigns one of three compute modes:

\begin{description}
  \item[Fast:] No thinking, low temperature (0.7), 8K max tokens. For greetings, simple factual queries.
  \item[Quality:] Thinking enabled (4{,}096 token budget), ABF at threshold 0.55, Engram active. Default mode.
  \item[Ultra:] DVTS tree search with $K{=}2$ candidates scored by ThinkPRM. For medical, complex reasoning.
\end{description}

\section{Web Search Pipeline}

The deep search pipeline (\texttt{deep\_search\_sota.py}, 834~lines) implements 8~stages:

\begin{enumerate}
  \item \textbf{Query expansion}: LLM generates 3--5 diversified sub-queries.
  \item \textbf{Parallel search}: Brave API + SearXNG simultaneously.
  \item \textbf{RRF fusion}: Reciprocal Rank Fusion merges result lists.
  \item \textbf{Deep fetch}: Trafilatura extracts clean text from URLs.
  \item \textbf{Chunking}: Split documents into overlapping chunks.
  \item \textbf{CRAG filtering}: Credibility-based chunk filtering.
  \item \textbf{Contradiction detection}: Flag conflicting claims across sources.
  \item \textbf{LLM synthesis}: Generate cited response from filtered chunks.
\end{enumerate}

\section{Self-Improvement Loop}

The system logs every interaction with quality scores from ThinkPRM-1.5B (step-level verification). Nightly processes run via systemd timers:

\begin{description}
  \item[03:00] LoRA training on quality-filtered pairs (score $\geq$ 4).
  \item[04:00] Engram WRITE: update domain tables from validated responses. Decay old entries ($>$30~days: halve weight; $>$90~days: delete).
  \item[Monday 02:00] DSPy MIPROv2 prompt optimization per domain.
  \item[Every 6h] ChromaDB RAG re-indexing.
\end{description}

As of March~29, 2026: 104~quality scores logged (mean: 3.04/5), 68~training pairs generated, 72~SPIN DPO pairs accumulated.

\section{Evaluation}

\subsection{Benchmark}

We evaluate on 10~questions across 4~domains: physiotherapy clinical reasoning (5), code generation (2), mathematics (1), tool calling (2). Questions are designed to test domain knowledge, reasoning, and system integration.

\begin{table}[h]
\centering
\caption{System benchmark progression.}
\begin{tabular}{lccccc}
\toprule
\textbf{Date} & \textbf{Kine} & \textbf{Code} & \textbf{Math} & \textbf{Tools} & \textbf{Total} \\
\midrule
Mar 25 (baseline) & 1/2 & 2/2 & 1/1 & 2/2 & 7/8 (88\%) \\
Mar 26 (ODO fixes) & 5/5 & 2/2 & 1/1 & 2/2 & 10/10 (100\%) \\
\bottomrule
\end{tabular}
\end{table}

The 50\%$\to$100\% improvement came from: (1)~fixing tool injection that suppressed web enrichment, (2)~reducing thinking budget from 4{,}096 to 2{,}048 tokens with ABF threshold 0.55 (eliminating timeouts on IQ3\_S), (3)~web search timeout increase from 30s to 120s.

\subsection{Limitations}

This benchmark has 10 questions --- a smoke test, not a rigorous evaluation. We do not claim MMLU or HumanEval scores. The self-improvement loop has been running for $<$2~weeks; longitudinal quality improvement is not yet demonstrated. The system is optimized for one hardware configuration (RTX~5060~Ti + DDR5).

\section{Negative Results and Lessons}

\begin{itemize}
  \item \textbf{Expert prefetch}: MLP predictor achieves 86.65\% hit@8 accuracy, but provides zero speedup because CPU GEMV for 430\,KB experts takes $<$50\,$\mu$s --- less than the callback overhead.
  \item \textbf{LoRA on 35B MoE}: does not fit 16\,GB even in 4-bit. Solutions identified: MeZO (zeroth-order optimizer), DeepSpeed ZeRO-2 CPU offload.
  \item \textbf{DFlash online}: $\tau{=}9.4$ offline but cannot deploy online due to GDN State Barrier (see companion paper).
  \item \textbf{Engram v1}: applying bias during thinking degrades performance by 8~pp. Response-only bias with lower $\alpha$ is critical.
  \item \textbf{NVFP4/MXFP4}: ineffective on 256 micro-experts (too small for format overhead).
\end{itemize}

\section{Related Work}

\paragraph{MoE inference on consumer hardware.}
KTransformers~(2024) addresses MoE offloading with CPU expert scheduling but provides no quality loop or orchestration.
PowerInfer~\cite{powerinfer2024} and PowerInfer-2 use neuron-level GPU/CPU splitting for MoE, targeting cold-start latency.
Ollama and LM~Studio provide easy deployment but no adaptive compute, no domain memory, no self-improvement.
vLLM and TensorRT-LLM achieve high throughput but require A100+ hardware.

\paragraph{Self-improving systems.}
Autoresearch~(Karpathy, 2025) proposes self-improving research agents but remains a concept paper.
Self-Play Fine-Tuning (SPIN)~\cite{spin2024} generates DPO pairs from the model's own outputs --- we use this for nightly training.
DSPy~\cite{dspy2024} optimizes prompts via Bayesian search --- we run MIPROv2 weekly.
STaR~(Zelikman et~al., 2022) bootstraps reasoning via self-taught rationalization.
Our contribution is integrating these into a deployed loop that runs unattended on consumer hardware.

\paragraph{Conditional memory.}
DeepSeek Engram~(arXiv:2601.07372) adds conditional memory to MoE with +3.4~MMLU, requiring end-to-end training.
We implement a lightweight inference-time approximation: N-gram hash lookup with logit biasing, achieving +3pp on our domain benchmark without any training.
RAG systems (HippoRAG~2, Self-RAG, CRAG) inject knowledge via context; Engram injects at the logit level --- complementary approaches.

\paragraph{Quantization.}
GPTQ, AWQ, and SpinQuant require calibration data and full model in memory.
Our RAMP pipeline~(companion paper) uses data-free NSDS sensitivity for consumer feasibility.

\paragraph{Speculative decoding.}
DFlash~(arXiv:2602.06036) introduces block diffusion drafting.
EAGLE and Medusa propose learned drafters.
We document that all these methods fail on hybrid GDN-attention architectures due to the State Barrier~(companion paper).

\section{Conclusion}

Chim\`ere demonstrates that a carefully integrated system can run a 35B MoE model at 93~tok/s on consumer hardware with self-improving quality. The system is fully open-source: runtime (\url{https://github.com/AIdevsmartdata/chimere}), orchestrator (\url{https://github.com/AIdevsmartdata/chimere-odo}), and quantization pipeline (\url{https://github.com/AIdevsmartdata/ramp-quant}).

\begin{thebibliography}{15}

\bibitem{dflash2024}
z-lab. DFlash: Accelerating autoregressive inference with block diffusion. \emph{arXiv:2602.06036}, 2024.

\bibitem{yang2024gated}
Yang, S., et~al. Gated delta networks: Improving Mamba-2 with delta rule. \emph{ICLR}, 2025.

\bibitem{engram2024}
DeepSeek-AI. Engram: Memory-augmented language models via conditional memory. \emph{arXiv:2601.07372}, 2024.

\bibitem{ktransformers2024}
KTransformers. CPU/GPU co-serving for MoE inference. GitHub, 2024.

\bibitem{powerinfer2024}
Song, Y., et~al. PowerInfer: Fast large language model serving with a consumer-grade GPU. \emph{MLSys}, 2024.

\bibitem{thinkprm2024}
Luo, L., et~al. Improve mathematical reasoning in language models by automated process supervision. \emph{arXiv:2406.06592}, 2024.

\bibitem{dspy2024}
Khattab, O., et~al. DSPy: Compiling declarative language model calls into self-improving pipelines. \emph{ICLR}, 2024.

\bibitem{spin2024}
Chen, Z., et~al. Self-play fine-tuning converts weak language models to strong language models. \emph{ICML}, 2024.

\bibitem{spinquant2024}
Liu, Z., et~al. SpinQuant: LLM quantization with learned rotations. \emph{arXiv:2405.16406}, 2024.

\end{thebibliography}

\end{document}
